\documentclass[aspectratio=169]{beamer}
\usetheme{Madrid}
\usepackage[spanish,es-nodecimaldot]{babel}
\usepackage[utf8]{inputenc}
\usepackage[T1]{fontenc}
\usepackage{tikz}
\usetikzlibrary{positioning,fit,backgrounds,arrows.meta}
\usepackage{xcolor}

% Paleta UNS + colores del mapa
\definecolor{azulUNS}{RGB}{30,60,120}
\usecolortheme[named=azulUNS]{structure}

\definecolor{tealBox}{RGB}{225,245,238}
\definecolor{tealBorde}{RGB}{15,110,86}
\definecolor{tealTxt}{RGB}{8,80,65}

\definecolor{coralBox}{RGB}{250,236,231}
\definecolor{coralBorde}{RGB}{153,60,29}
\definecolor{coralTxt}{RGB}{113,43,19}

\definecolor{purpBox}{RGB}{238,237,254}
\definecolor{purpBorde}{RGB}{83,74,183}
\definecolor{purpTxt}{RGB}{38,33,92}

\definecolor{azulBox}{RGB}{230,241,251}
\definecolor{azulBorde}{RGB}{24,95,165}

\definecolor{grisBox}{RGB}{241,239,232}
\definecolor{grisBorde}{RGB}{95,94,90}

\title[MAD1 -- U6]{Métodos de Análisis de Datos I}
\subtitle{Unidad 6 -- Bloque 1: Otros paradigmas de aprendizaje}
\author{Departamento de Matemática -- UNS}
\date{2026}

\begin{document}

\begin{frame}{El mapa del aprendizaje automático}
\begin{center}
\begin{tikzpicture}[
  font=\scriptsize,
  zona/.style={
    rectangle, rounded corners=3pt, draw=tealBorde, line width=0.4pt,
    fill=tealBox, text=tealTxt, align=left, inner sep=4pt,
    text width=3.6cm, minimum height=0.85cm},
  chip/.style={
    rectangle, rounded corners=4pt, draw=purpBorde, line width=0.4pt,
    fill=purpBox, text=purpTxt, align=center, inner sep=3pt,
    text width=2.1cm, minimum height=0.45cm},
]

% --- Region izquierda: aprende de datos ---
\node[zona] (sup)  at (0,0)    {\textbf{Supervisado}\\\tiny etiquetas humanas completas};
\node[zona] (semi) at (0,-1.1) {\textbf{Semi-supervisado}\\\tiny muy pocas etiquetas};
\node[zona] (auto) at (0,-2.2) {\textbf{Auto-supervisado}\\\tiny etiquetas sintéticas};
\node[zona] (nosup)at (0,-3.3) {\textbf{No supervisado}\\\tiny sin etiquetas};

% Eje de etiquetas a la izquierda
\draw[-{Latex[length=2mm]}, azulBorde, line width=0.6pt]
  ([xshift=-3pt]sup.north west) -- ([xshift=-3pt]nosup.south west)
  node[midway, left, rotate=90, anchor=south, text=azulBorde, font=\tiny]
  {etiquetas humanas: menos $\leftarrow$ más};

\begin{scope}[on background layer]
  \node[rectangle, rounded corners=6pt, draw=azulBorde, line width=0.4pt,
        fill=azulBox, fit=(sup)(semi)(auto)(nosup),
        inner sep=10pt, label={[anchor=north west, text=azulBorde,
        font=\scriptsize\bfseries]north west:Aprende de un conjunto de datos}] (regdatos) {};
\end{scope}

% --- Region derecha: refuerzo ---
\node[rectangle, rounded corners=6pt, draw=coralBorde, line width=0.4pt,
      fill=coralBox, text=coralTxt, align=left, inner sep=10pt,
      text width=3.4cm, minimum height=4.3cm, anchor=north]
      (rl) at (6.6,1.15)
      {\textbf{Por refuerzo (RL)}\\[4pt]
       \tiny Aprende actuando, no de un dataset fijo\\[5pt]
       \tiny Señal: recompensa por cada acción\\[5pt]
       \tiny El entorno responde a lo que el agente hace\\[5pt]
       \tiny Ej.: AlphaGo, robótica, control de tráfico};

% --- Franja inferior: transversales ---
\node[chip] (c1) at (-1.7,-4.7) {Transfer learning};
\node[chip] (c2) at (0.7,-4.7)  {Federado};
\node[chip] (c3) at (3.1,-4.7)  {Activo};
\node[chip] (c4) at (5.5,-4.7)  {Few / zero-shot};

\begin{scope}[on background layer]
  \node[rectangle, rounded corners=6pt, draw=grisBorde, line width=0.4pt,
        fill=grisBox, fit=(c1)(c2)(c3)(c4), inner sep=9pt,
        label={[anchor=north west, font=\scriptsize\bfseries]north west:%
        Enfoques transversales \textmd{\scriptsize(se combinan con cualquiera de los anteriores)}}] {};
\end{scope}

\end{tikzpicture}
\end{center}

\vspace{0.2cm}
{\scriptsize El eje que separa a los cuatro primeros es \textbf{cuántas etiquetas humanas} hay. RL aprende por interacción, no de un dataset fijo. Los enfoques transversales se montan sobre cualquiera (p.\,ej.\ GPT = auto-supervisado + transfer).}

\end{frame}

\end{document}
